\documentclass[11pt]{article}

\usepackage[a4paper,margin=27mm]{geometry}
\usepackage[T1]{fontenc}
\usepackage[utf8]{inputenc}
\usepackage{lmodern}
\usepackage{microtype}
\usepackage{amsmath,amssymb,amsfonts,amsthm,mathtools,bm}
\usepackage{booktabs}
\usepackage{xcolor}
\usepackage{enumitem}
\usepackage{natbib}
\usepackage{hyperref}
\usepackage[nameinlink,capitalize,noabbrev]{cleveref}

\hypersetup{
  colorlinks=true,
  linkcolor=blue!50!black,
  citecolor=blue!50!black,
  urlcolor=blue!50!black,
  pdftitle={Universal Tangent-Kernel Closure in Two-Factor Linear Networks, and Its Failure Beyond},
  pdfauthor={Redha Moulla}
}

\setlist[itemize]{leftmargin=1.5em,itemsep=0.25em,topsep=0.4em}
\setlist[enumerate]{leftmargin=1.7em,itemsep=0.25em,topsep=0.4em}
\allowdisplaybreaks

\newtheorem{theorem}{Theorem}[section]
\newtheorem{proposition}[theorem]{Proposition}
\newtheorem{lemma}[theorem]{Lemma}
\newtheorem{corollary}[theorem]{Corollary}
\newtheorem{definition}[theorem]{Definition}
\newtheorem{remark}[theorem]{Remark}
\newtheorem*{remark*}{Remark}

\crefname{theorem}{Theorem}{Theorems}
\Crefname{theorem}{Theorem}{Theorems}
\crefname{proposition}{Proposition}{Propositions}
\Crefname{proposition}{Proposition}{Propositions}
\crefname{lemma}{Lemma}{Lemmas}
\Crefname{lemma}{Lemma}{Lemmas}
\crefname{corollary}{Corollary}{Corollaries}
\Crefname{corollary}{Corollary}{Corollaries}
\crefname{definition}{Definition}{Definitions}
\Crefname{definition}{Definition}{Definitions}
\crefname{remark}{Remark}{Remarks}
\Crefname{remark}{Remark}{Remarks}

\newcommand{\R}{\mathbb{R}}
\newcommand{\Spos}{\mathbb{S}_{+}}
\newcommand{\norm}[1]{\left\lVert #1\right\rVert}
\newcommand{\dd}{\mathrm{d}}
\newcommand{\Ker}{\operatorname{ker}}
\newcommand{\rank}{\operatorname{rank}}

\title{\bfseries Universal Tangent-Kernel Closure in Two-Factor Linear Networks,\\
and Its Failure Beyond}
\author{Redha Moulla\\[2mm]\small AXIA, France}
\date{August 2026}

\begin{document}
\maketitle

\begin{abstract}
Finite-width gradient flow in a two-factor linear network admits an exact projection from factor space to the joint prediction--kernel state $(f,K)$.  For every scalar-output bias-free model
\[
  f=XW^\top a,
  \qquad G=XX^\top,
\]
with arbitrary hidden width, input dimension, finite dataset, and differentiable loss, the answer is yes:
\[
  K=XW^\top WX^\top+\norm{a}^2G,
  \qquad
  \dot K=-(Gr)f^\top-f(Gr)^\top-2(f^\top r)G.
\]
The contribution is the closure statement and its consequences: the parameter representation disappears globally, the law is independent of the nontrivial balancing leaf $WW^\top-aa^\top$, and the kernel equation is driven by $(f,r)$ rather than by $K$ itself.  Eliminating $K$ gives an exact prediction-only Volterra equation with full, non-fading history and no remainder once $K_0$ is fixed.  Thus the same dynamics are Markovian in $(f,K)$ and memory-bearing in $f$ alone.

For scalar deep-linear networks, we derive the finite hierarchy
\[
  \dot e_k=-2fr(d-k+1)e_{k-1},
  \qquad K=e_{d-1},
  \qquad \dot K=-4fr e_{d-2}.
\]
At two factors, $e_{d-2}=1$ and the universal closure holds.  From three factors onward, regular fibers of $(e_d,e_{d-1})$ contain curves along which $e_{d-2}$ varies, so no leaf-independent law $\dot K=\mathcal F(f,K,r)$ exists on the full positive state space.  Conditioning on the conserved balancing leaf restores exact closure; at depth three, $(f,K,I)$ is simply the complete symmetric state in invariant coordinates.  The result is therefore a precise transition from universal driven closure to leaf-conditioned closure.  It also separates memory created by eliminating an observable from the distinct closure problem posed by nonlinear networks.
\end{abstract}

\section{Introduction}
\label{sec:introduction}

Let a scalar-output model be evaluated on a finite dataset $\{x_i\}_{i=1}^n$.  Write
\begin{equation}
  f(w)=\bigl(f(x_1;w),\ldots,f(x_n;w)\bigr)^\top\in\R^n,
  \qquad
  J(w)=D_w f(w)\in\R^{n\times p},
\end{equation}
and define the empirical tangent kernel
\begin{equation}
  K(w)=J(w)J(w)^\top\in\Spos^n.
\end{equation}
For a differentiable loss $L(w)=\ell(f(w))$, let $r(f)=\nabla_f\ell(f)$.  Euclidean gradient flow gives
\begin{equation}
  \dot w=-J^\top r,
  \qquad
  \dot f=-Kr,
  \qquad
  \dot L=-r^\top Kr\leq0.
  \label{eq:basic-flow}
\end{equation}
At infinite width, the neural tangent kernel may become asymptotically constant, yielding a closed function-space description \citep{jacot2018ntk,lee2019linearized}.  At finite width, differentiating $K$ generally introduces higher derivatives and leads to the neural tangent hierarchy \citep{huang2020nth}.

The question here is exact and finite-dimensional.  On a parameter domain $\mathcal D$, call $(f,K)$ a \emph{closed state} if there exists a single-valued map $\mathcal F$ such that
\begin{equation}
  \dot K(w)=\mathcal F\bigl(f(w),K(w),r(f(w))\bigr)
  \qquad \text{for every }w\in\mathcal D.
  \label{eq:closure-definition}
\end{equation}
Equivalently, the parameter-space vector field for $\dot K$ factors through the observation map $w\mapsto(f,K)$.  This is stronger than having an equation along one trajectory and weaker than demanding an autonomous field on kernel space alone.

The main positive result is an exact closure for every bias-free two-factor linear network with scalar output.  It is universal with respect to hidden width, input dimension, finite dataset, differentiable loss, and parameter state.  It is also structurally degenerate: $K$ drives $f$ through \eqref{eq:basic-flow}, but $\dot K$ has no direct dependence on $K$.  Once the prediction path is known, the kernel follows by quadrature.  This one-sided structure permits exact elimination of $K$, producing a prediction-only equation with explicit memory.

The scalar deep-linear hierarchy then identifies the first failure of this property.  Two factors close because the hierarchy terminates at the constant $e_0=1$.  At three factors, the kernel velocity depends on $e_1$, which varies along regular fibers with fixed output and fixed tangent kernel.  The missing information is not an active hidden state in this solvable model: squared-layer differences are conserved, and conditioning on the balancing leaf restores a Markovian closure.

\paragraph{Relation to prior work and contribution.}
The decomposition
\[
  K(x,x')=x^\top W^\top W x'+\norm{a}^2x^\top x'
\]
appears explicitly in analyses of two-layer linear alignment, including the silent-alignment study of \citet{atanasov2022silent}.  Earlier work develops exact and asymptotic dynamics for two-factor and deep-linear networks, balancing laws, and structured mode-wise solutions \citep{saxe2014deep,arora2018optimization,du2018balanced,xu2025three}.  This note places these ingredients in an exact state-reduction framework and establishes four linked results:
\begin{enumerate}
  \item the full parameter vector field projects globally to $(f,K)$ for arbitrary finite data and differentiable loss, without whitening, teacher, small-initialization, or mode-decoupling assumptions;
  \item the projected law is independent of the nontrivial balancing leaf, yet is driven rather than kernel-autonomous;
  \item exact elimination of the driven variable produces a non-fading Volterra representation of the same dynamics;
  \item the first scalar deep-linear failure is characterized directly by the geometry of the fibers, and exact closure is recovered after conditioning on conserved leaf coordinates.
\end{enumerate}
Together, these results identify a global projectability law at two factors and locate the first scalar factorization at which conserved leaf information becomes dynamically necessary.


\section{Universal driven closure for two factors}
\label{sec:two-factor}

Let $X\in\R^{n\times d_x}$ contain the data points as rows and consider
\begin{equation}
  f=XW^\top a,
  \qquad
  W\in\R^{m\times d_x},
  \qquad
  a\in\R^m.
  \label{eq:two-factor-model}
\end{equation}
Let $G=XX^\top$ be the fixed data Gram matrix and $r=\nabla_f\ell(f)$.

\begin{proposition}[Universal two-factor tangent-kernel closure]
\label[proposition]{prop:two-factor-closure}
For every hidden width $m$, input dimension $d_x$, finite dataset, differentiable loss, and parameter state,
\begin{equation}
  K=XW^\top WX^\top+\norm{a}^2G,
  \label{eq:two-factor-kernel}
\end{equation}
and gradient flow satisfies
\begin{equation}
  \boxed{
  \dot f=-Kr,
  \qquad
  \dot K=-(Gr)f^\top-f(Gr)^\top-2(f^\top r)G.}
  \label{eq:two-factor-kdot}
\end{equation}
Hence $(f,K)$ is an exact closed state.  The $K$-component depends only on $(f,r)$ and the fixed data geometry $G$, not on the hidden factors or directly on $K$.
\end{proposition}

\begin{proof}
For sample $i$,
\begin{equation}
  \nabla_W f_i=a x_i^\top,
  \qquad
  \nabla_a f_i=Wx_i,
\end{equation}
which gives \eqref{eq:two-factor-kernel}.  Put $s=X^\top r$.  Gradient flow is
\begin{equation}
  \dot W=-a s^\top,
  \qquad
  \dot a=-Ws.
  \label{eq:two-factor-parameter-flow}
\end{equation}
Because $G=XX^\top$ is fixed, $\dot G=0$.  Differentiating \eqref{eq:two-factor-kernel} gives
\begin{equation}
  \dot K
  =X(\dot W^\top W+W^\top\dot W)X^\top
   +\left(\frac{\dd}{\dd t}\norm{a}^2\right)G.
  \label{eq:two-factor-proof-step-one}
\end{equation}
Here the scalar derivative multiplying the fixed matrix $G$ is
\[
  \frac{\dd}{\dd t}\norm{a}^2
  =2\langle a,\dot a\rangle
  =-2a^\top Ws
  =-2f^\top r.
\]
Substituting \eqref{eq:two-factor-parameter-flow} into \eqref{eq:two-factor-proof-step-one} therefore yields
\begin{equation}
  \dot K
  =-(Xs)(a^\top WX^\top)-(XW^\top a)(Xs)^\top
   -2(a^\top Ws)G.
  \label{eq:two-factor-proof-substitution}
\end{equation}
Thus the last contribution is $\bigl(2\langle a,\dot a\rangle\bigr)G=-2(f^\top r)G$; it comes from differentiating $\norm{a}^2$, never from differentiating the fixed matrix $G$.  Since $Xs=Gr$ and $XW^\top a=f$, \eqref{eq:two-factor-kdot} follows.
\end{proof}

\subsection{Driven closure rather than a constitutive self-law}

Define
\begin{equation}
  B_G(f,r)=-(Gr)f^\top-f(Gr)^\top-2(f^\top r)G.
  \label{eq:BG}
\end{equation}
The exact reduced system is
\begin{equation}
  \dot f=-Kr(f),
  \qquad
  \dot K=B_G\bigl(f,r(f)\bigr).
  \label{eq:joint-driven-system}
\end{equation}
It is autonomous on the product state $(f,K)$, but it is not a constitutive field $K\mapsto\dot K$ on the positive-semidefinite cone.  The kernel has no instantaneous self-dynamics: its influence on its own future is indirect, through the prediction path that it helps generate.  In this precise sense the closure is \emph{driven}.  Consequently, the exact case in which the hidden representation disappears is also a case in which a kernel-only gradient-flow question becomes vacuous unless one freezes the prediction, treats it as a control, or works on the product space.

\subsection{A nontrivial foliation that the closure does not need}

\begin{proposition}[Balancing invariant and leaf independence]
\label[proposition]{prop:two-factor-invariant}
The matrix
\begin{equation}
  C=WW^\top-aa^\top
  \label{eq:two-factor-invariant}
\end{equation}
is conserved under \eqref{eq:two-factor-parameter-flow}.  Nevertheless, the closed law \eqref{eq:two-factor-kdot} is independent of $C$.
\end{proposition}

\begin{proof}
Using \eqref{eq:two-factor-parameter-flow},
\begin{align}
  \frac{\dd}{\dd t}(WW^\top)
  &=-a s^\top W^\top-Ws a^\top,\\
  \frac{\dd}{\dd t}(aa^\top)
  &=-Ws a^\top-a s^\top W^\top.
\end{align}
The derivatives are equal, so $\dot C=0$.  The formula \eqref{eq:two-factor-kdot} contains no $C$, proving leaf independence.
\end{proof}

This separates two notions that are often conflated.  An invariant foliation may exist without its labels being required by a reduced state.  The relevant question is not whether invariants exist, but whether the projected velocity changes across the leaves identified by the chosen observable.

\subsection{Exact elimination and a non-fading memory law}

The $K$-equation in \eqref{eq:joint-driven-system} is integrable once the prediction path is known.  Define
\begin{equation}
  \mathcal Q_G(f,u)
  =(Gu)f^\top+f(Gu)^\top+2(f^\top u)G
  =-B_G(f,u).
  \label{eq:QG}
\end{equation}

\begin{proposition}[Exact prediction-only Volterra equation]
\label[proposition]{prop:volterra}
Fix $K_0$ and write $r_t=r(f_t)$.  The Markovian system \eqref{eq:joint-driven-system} is equivalent to
\begin{equation}
  K_t=K_0-\int_0^t\mathcal Q_G(f_s,r_s)\,\dd s
  \label{eq:K-integral}
\end{equation}
and the closed prediction-only equation
\begin{equation}
  \boxed{
  \dot f_t=-K_0r_t
  +\int_0^t\mathcal Q_G(f_s,r_s)r_t\,\dd s.}
  \label{eq:volterra}
\end{equation}
Equivalently,
\begin{align}
  \dot f_t={}&-K_0r_t
  +\int_0^t\Bigl[
    (Gr_s)(f_s^\top r_t)
    +f_s(r_s^\top G r_t)\nonumber\\
  &\hspace{31mm}+2(f_s^\top r_s)Gr_t
  \Bigr] \,\dd s.
  \label{eq:volterra-expanded}
\end{align}
\end{proposition}

\begin{proof}
Integrating $\dot K=-\mathcal Q_G(f,r)$ gives \eqref{eq:K-integral}.  Substitution into $\dot f_t=-K_t r_t$ gives \eqref{eq:volterra}.  Conversely, reconstructing $K$ by \eqref{eq:K-integral} recovers both equations in \eqref{eq:joint-driven-system}.
\end{proof}

Equation \eqref{eq:volterra} contains the complete past $s\in[0,t]$ without an explicit lag-decay factor.  Old contributions may become small because the states $(f_s,r_s)$ become small, but the representation has no structural fading kernel that would justify truncating the history to a short window.  Conditional on $K_0$, there is no additional stochastic or orthogonal-dynamics remainder: the integral carries the entire effect of the eliminated kernel.  The same deterministic training dynamics are therefore Markovian in $(f,K)$ and history-dependent in $f$ alone.  Memory is a property of the chosen observable, not of the underlying vector field in isolation.  This is an elementary, completely resolved instance of the projection principle emphasized by Mori--Zwanzig theory \citep{zwanzig1961memory,mori1965transport,chorin2002memory}.


\section{Scalar deep-linear hierarchy and failure beyond two factors}
\label{sec:deep-linear}

Consider the scalar width-one depth-$d$ linear network
\begin{equation}
  f(w)=\prod_{i=1}^d w_i
  \label{eq:scalar-network}
\end{equation}
with scalar residual $r=\ell'(f)$.  On a fixed sign component, put
\begin{equation}
  q_i=w_i^2>0,
  \qquad
  e_k(q)=\sum_{1\leq i_1<\cdots<i_k\leq d}q_{i_1}\cdots q_{i_k},
  \qquad e_0=1.
\end{equation}
Then $f^2=e_d$ and
\begin{equation}
  K=\sum_{i=1}^d\left(\frac{\partial f}{\partial w_i}\right)^2
  =e_{d-1}(q).
  \label{eq:scalar-K}
\end{equation}
Gradient flow gives the common translation law
\begin{equation}
  \dot q_i=-2fr
  \qquad \text{for every }i.
  \label{eq:q-common-velocity}
\end{equation}

\begin{proposition}[Finite symmetric-polynomial hierarchy]
\label[proposition]{prop:symmetric-hierarchy}
For $k=1,\ldots,d$,
\begin{equation}
  \boxed{\dot e_k=-2fr(d-k+1)e_{k-1}.}
  \label{eq:hierarchy}
\end{equation}
In particular,
\begin{equation}
  \dot f=-e_{d-1}r=-Kr,
  \qquad
  \boxed{\dot K=-4fr e_{d-2}.}
  \label{eq:K-hierarchy}
\end{equation}
Thus $(f,e_1,\ldots,e_{d-1})$ is an exact finite-dimensional state.
\end{proposition}

\begin{proof}
Since every $q_i$ has velocity $-2fr$,
\begin{equation}
  \dot e_k=-2fr\sum_{i=1}^d\frac{\partial e_k}{\partial q_i}.
\end{equation}
Every monomial of degree $k-1$ occurs in exactly $d-k+1$ of the derivatives, so
\begin{equation}
  \sum_i\frac{\partial e_k}{\partial q_i}=(d-k+1)e_{k-1}.
\end{equation}
The stated equations follow.
\end{proof}

At $d=1$, $K\equiv1$.  At $d=2$, \eqref{eq:K-hierarchy} gives
\begin{equation}
  \dot f=-Kr,
  \qquad
  \dot K=-4fr,
  \label{eq:scalar-depth-two}
\end{equation}
which is the scalar specialization of \cref{prop:two-factor-closure}.  From $d=3$ onward, $e_{d-2}$ is a new coordinate.  The next theorem shows directly that it is not determined by $(f,K)$ on generic fibers.

\begin{theorem}[Regular-fiber obstruction from three factors onward]
\label[theorem]{thm:fiber-obstruction}
Let $d\geq3$ and define
\begin{equation}
  \pi_d(q)=\bigl(e_d(q),e_{d-1}(q)\bigr),
  \qquad q\in\R_{>0}^d.
\end{equation}
At every point for which the reciprocals $u_i=1/q_i$ take at least three distinct values,
\begin{equation}
  \rank D\bigl(e_d,e_{d-1},e_{d-2}\bigr)=3.
  \label{eq:rank-three}
\end{equation}
Consequently, the fiber of $\pi_d$ is locally a smooth manifold of dimension $d-2$ containing curves along which $e_{d-2}$ is nonconstant.  For $d=3$ the fiber itself is locally a curve, and $e_1$ varies along it.

Therefore, for any loss and output value with $fr\neq0$, there is no leaf-independent single-valued law
\begin{equation}
  \dot K=\mathcal F(f,K,r)
  \label{eq:no-closure}
\end{equation}
on the full positive state space.  Three factors are the first scalar linear architecture at which the universal two-factor closure fails.
\end{theorem}

\begin{proof}
The change of variables $u_i=1/q_i$ is a diffeomorphism of the positive orthant.  Put
\begin{equation}
  P=e_d(q),
  \qquad
  S=\frac{e_{d-1}(q)}{e_d(q)}=\sum_i u_i,
  \qquad
  T=\frac{e_{d-2}(q)}{e_d(q)}=\sum_{i<j}u_i u_j.
  \label{eq:reciprocal-coordinates}
\end{equation}
The transformations $(e_d,e_{d-1},e_{d-2})\leftrightarrow(\log P,S,T)$ are locally invertible for $P>0$, so it suffices to study the differentials in $u$-coordinates:
\begin{equation}
  \dd\log P=-\sum_i\frac{\dd u_i}{u_i},
  \qquad
  \dd S=\sum_i\dd u_i,
  \qquad
  \dd T=\sum_i(S-u_i)\dd u_i.
  \label{eq:three-differentials}
\end{equation}
The first two forms are independent unless all $u_i$ are equal.  Suppose that $\dd T=A\,\dd S+B\,\dd\log P$.  Comparing the coefficient of $\dd u_i$ gives
\begin{equation}
  S-u_i=A-\frac{B}{u_i}.
\end{equation}
Hence every $u_i$ is a root of the same quadratic
\begin{equation}
  z^2-(S-A)z-B=0.
\end{equation}
This is impossible when at least three distinct reciprocal values occur.  The three differentials are therefore independent, proving \eqref{eq:rank-three}.

The constant-rank theorem now makes the level set of $(e_d,e_{d-1})$ a local manifold of dimension $d-2$.  Since $\dd e_{d-2}$ is not in the span of $\dd e_d$ and $\dd e_{d-1}$, it has a nonzero derivative in some tangent direction to that fiber; a smooth curve tangent to that direction keeps $(e_d,e_{d-1})$ fixed while changing $e_{d-2}$.  Along the positive sign component, fixed $e_d$ means fixed $f$, and fixed $e_{d-1}$ means fixed $K$.  Equation \eqref{eq:K-hierarchy} then gives different values of $\dot K$ whenever $fr\neq0$.
\end{proof}

\begin{remark*}[Explicit depth-three witness]
For $d=3$, take
\[
  q^A=\left(\frac25,\frac54,2\right),
  \qquad
  q^B=\left(\frac59,\frac35,3\right).
\]
Then $e_3(q^A)=e_3(q^B)=1$ and $e_2(q^A)=e_2(q^B)=19/5$, whereas
\[
  e_1(q^A)=\frac{73}{20},
  \qquad
  e_1(q^B)=\frac{187}{45}.
\]
Thus the two states have the same positive output $f=1$ and the same tangent kernel $K=19/5$, but \eqref{eq:K-hierarchy} gives different kernel velocities whenever $r(1)\neq0$.
\end{remark*}

The condition of three distinct reciprocal coordinates is open and dense.  The obstruction is therefore not an isolated pair of parameter states but a generic local property of the positive deep-linear state space.  It rules out a \emph{universal} closure on that space.  It does not rule out closure on a selected invariant leaf or on a more restrictive set of trajectories.

\subsection{Exact closure on each balancing leaf}

Equation \eqref{eq:q-common-velocity} implies that
\begin{equation}
  c_i=q_i-q_d,
  \qquad i=1,\ldots,d-1,
  \qquad \dot c_i=0.
  \label{eq:leaf-labels}
\end{equation}
Fix a positive leaf $c$ and write $q_d=s$, $q_i=s+c_i$.  Define
\begin{equation}
  \kappa_c(s)=e_{d-1}(s+c_1,\ldots,s+c_{d-1},s).
\end{equation}

\begin{proposition}[Leaf-conditioned exact closure]
\label[proposition]{prop:leaf-closure}
On every interval where all $q_i$ are positive,
\begin{equation}
  \kappa_c'(s)=2e_{d-2}(q(s,c))>0.
  \label{eq:kappa-derivative}
\end{equation}
Hence $s=\kappa_c^{-1}(K)$ and
\begin{equation}
  \boxed{
  \dot K=-4fr\,e_{d-2}\bigl(q(\kappa_c^{-1}(K),c)\bigr),
  \qquad \dot c=0.}
  \label{eq:leaf-closure}
\end{equation}
Thus the failure in \cref{thm:fiber-obstruction} is transverse to the balancing foliation: a conserved label, not temporal history, restores exact closure.
\end{proposition}

\begin{proof}
Differentiation under a common translation and \cref{prop:symmetric-hierarchy} give
\begin{equation}
  \frac{\dd}{\dd s}e_{d-1}(q(s,c))=2e_{d-2}(q(s,c)).
\end{equation}
Positivity makes this derivative strictly positive.  Inversion and \eqref{eq:K-hierarchy} yield \eqref{eq:leaf-closure}.
\end{proof}

\subsection{Depth three in invariant coordinates}

For $d=3$, let
\begin{equation}
  I(q)=\frac12\sum_{i=1}^3(q_i-\bar q)^2,
  \qquad
  \bar q=\frac13e_1(q).
  \label{eq:I}
\end{equation}
The common translation \eqref{eq:q-common-velocity} gives $\dot I=0$, and
\begin{equation}
  I=\frac{e_1^2}{3}-e_2.
  \label{eq:I-identity}
\end{equation}
Since $K=e_2$ and $q_i>0$,
\begin{equation}
  e_1=\sqrt{3(K+I)}.
\end{equation}
The exact closed system is therefore
\begin{equation}
  \boxed{
  \dot f=-Kr,
  \qquad
  \dot K=-4fr\sqrt{3(K+I)},
  \qquad
  \dot I=0.}
  \label{eq:depth-three-closure}
\end{equation}
Because $f^2=e_3$, $K=e_2$, and $I$ recovers $e_1$, the state $(f,K,I)$ is simply the complete symmetric state $(e_3,e_2,e_1)$ written in invariant coordinates.  The one-scalar augmentation is not a further compression of the hierarchy.


\input{sections/04_interpretation_scope.tex}

\section{Conclusion}

Two-factor linear networks provide an exact finite-width closure that is broader than the special regimes in which their learning dynamics are usually solved.  For arbitrary finite data, differentiable loss, hidden width, and input dimension, the parameter flow projects to $(f,K)$, independently of the balancing leaf.  The resulting law is driven rather than kernel-autonomous, and exact elimination of $K$ turns it into a full-history Volterra equation for $f$.  The same dynamics are therefore Markovian or memory-bearing according to the observable retained.

The scalar deep-linear hierarchy pinpoints the first rupture.  From three factors onward, generic fibers with fixed output and fixed tangent kernel carry varying $e_{d-2}$ and hence varying kernel velocity.  This failure is not irreducible memory: balancing labels restore closure, and at depth three the required invariant coordinate completes the symmetric hierarchy.  The transition is thus from universal leaf-independent closure to leaf-conditioned closure.  Whether analogous reductions hold on training-reachable sets of nonlinear networks is a separate empirical question.


\bibliographystyle{plainnat}
\bibliography{references}

\end{document}
